% =============================================================================
% CAPITULO 6 - RESULTADOS
% =============================================================================
\chapter{Resultados}
\label{ch:resultados}

Este capítulo reporta los resultados del \CCV sobre el Mundial Qatar
2022 en dos bloques. El primero, validación interna, cubre la
calibración de los modelos predictivos del \textit{pipeline} y la
robustez del identificador causal. El segundo, resultados sustantivos,
descompone el efecto del \textit{shock} a nivel poblacional, lo
contrasta con la heterogeneidad individual estimada por el CATE
jerárquico bayesiano y aterriza en los jugadores con respuesta
significativa por contexto.


\section{Calibración de los modelos predictivos}
\label{sec:res-calibracion}

La capa predictiva sobre la que descansa el cuasi-experimento de
\textit{near-miss} es el modelo Post-Shot~$xG$. La figura
\ref{fig:psxg-calibration} reporta su calibración sobre la fase de
grupos del Mundial y sobre el resto de torneos de entrenamiento.

\begin{figure}[h]
  \centering
  \includegraphics[width=\textwidth]{fig_cap6_psxg_calibration}
  \caption{Curva de fiabilidad y descomposición de Brier del modelo PSxG.}
  \label{fig:psxg-calibration}
\end{figure}

Las dos curvas se mantienen sobre la diagonal de calibración perfecta
en todo el rango, con desviaciones inferiores a un punto porcentual
en cada uno de los cuatro \textit{bins} de la regresión isotónica.
La AUC \textit{out-of-fold} (OOF) es 0{,}968 y la AUC sobre el
\textit{holdout} del Mundial es 0{,}976. La AUC alta del modelo
respecto al rango clásico de los $xG$ pre-disparo
\parencite{AnzerBauer2021} es consistente con la ventaja informativa del
PSxG, que incorpora \textit{end\_y} y \textit{end\_z} del disparo y
la configuración 360 \textit{freeze-frame}
en el instante del lanzamiento; el detalle del análisis y la
comparación contra la literatura PSxG están en el
anexo~\ref{anexo:val-audit-models}. El
\textit{Expected Calibration Error} (ECE) sobre el \textit{holdout}
es 0{,}011, y el componente de \textit{reliability} en la
descomposición de Brier es 0{,}0002 ---dos órdenes de magnitud por
debajo del baseline $xG$ pre-disparo de StatsBomb tomado como
referencia, con la salvedad de que ambos modelos operan sobre tareas
distintas (PSxG post-disparo frente a $xG$ pre-disparo). La componente
de \textit{resolution} es 0{,}0713 sobre el \textit{holdout},
aproximadamente tres veces y media superior al mismo baseline. La
combinación de reliability bajo y resolution alto respalda el uso del
modelo como ancla técnica del corpus de \textit{near-miss} sin sesgar
la asignación entre tratamiento y control.

La capa de \textit{Win Probability} (WP) sigue el mismo protocolo:
ajuste sobre Wyscout 2017/18 más StatsBomb fuera del Mundial y
recalibración por \textit{temperature scaling} sobre la fase de
grupos. La curva de calibración resultante sobre la fase de
grupos del Mundial se mantiene próxima a la diagonal de calibración
perfecta, con un parámetro de temperatura aprendido ligeramente
inferior a la unidad ($T \approx 0{,}95$), consistente con un
desplazamiento de distribución modesto desde la liga europea al
torneo de selecciones. La fase
eliminatoria queda como \textit{holdout} final donde no se ajusta
ningún parámetro.


\section{Identificación causal: pre-tendencias y sensibilidad}
\label{sec:res-identificacion}

La validación del brazo causal DiD se apoya en tres comprobaciones
independientes sobre los coeficientes \textit{within-player}. La
primera es el test $F$ de pre-tendencias agregadas
\parencite{Roth2022Pretest}, que no rechaza la hipótesis nula de
paralelismo pre-tratamiento en ninguno de los ocho contrastes canal
$\times$ contexto al nivel del cinco por ciento. Los coeficientes
pre-shock son indistinguibles de cero en los diez minutos previos al
gol en los cuatro canales, lo que es consistente con la asunción de
tendencias paralelas necesaria para la identificación del efecto del
\textit{shock} ---reconociendo que la potencia del test es limitada
con el tamaño de muestra disponible, motivo por el que las cotas
\textit{HonestDiD} se reportan como verificación complementaria.

La segunda comprobación es la sensibilidad temporal de los
coeficientes a la elección de la ventana pre/post. La figura
\ref{fig:window-sensitivity} reporta el efecto medio estimado sobre
horizontes de $\pm 3$, $\pm 5$, $\pm 7$, $\pm 10$ y $\pm 15$ minutos
por canal y por dirección del \textit{shock}.

\begin{figure}[h]
  \centering
  \includegraphics[width=\textwidth]{fig_cap6_window_sensitivity}
  \caption{Sensibilidad del ATE al horizonte pre/post.}
  \label{fig:window-sensitivity}
\end{figure}

Los coeficientes ofensivo y \textit{off-ball} son estables a través
de los cinco horizontes y cercanos a cero en todos ellos. El
coeficiente defensivo también es estable, con dispersión adicional
en los horizontes más cortos por la menor muestra efectiva. El
coeficiente físico es el más sensible al horizonte: en la ventana
estándar de $\pm 10$\,minutos los efectos son indistinguibles de
cero, pero al extender a $\pm 15$\,minutos aparecen efectos
estadísticamente significativos al cinco por ciento ---positivos
post-marcar ($+0{,}053$\,desviaciones estándar, $p = 0{,}0002$) y
modestos post-encajar ($+0{,}032$, $p = 0{,}047$). El patrón es
consistente con una respuesta física diferida que no se manifiesta
en la ventana inmediata pero sí en una ventana extendida; la
elección de $\pm 10$\,minutos como horizonte de referencia es
deliberada y se discute en el capítulo siguiente.

La tercera comprobación es la sensibilidad a violaciones de
tendencias paralelas mediante las cotas \textit{HonestDiD}
\parencite{RambachanRoth2023}, reportadas en el
anexo~\ref{anexo:validaciones}. Para todos los canales y direcciones,
las cotas robustecidas correspondientes a los tres niveles
$M \in \{0{,}5,\, 1,\, 2\}$ contienen el cero, lo que indica que los
coeficientes nulos del ATE poblacional son estables incluso si la
asunción de tendencias paralelas se relaja hasta el doble de la
máxima desviación pre-tratamiento observada.


\section{Efecto poblacional medio del shock}
\label{sec:res-ate-poblacional}

Antes de reportar los coeficientes causales, la
tabla~\ref{tab:corpus-shocks} describe la composición del corpus de
\textit{shocks} sobre el que se estiman. El torneo aporta 172
\textit{shocks} de gol únicos repartidos en 64 partidos, con sesgo
hacia la fase de grupos por volumen ---más partidos por equipo--- y
hacia los minutos centrales y finales de cada parte. La mitad de los
\textit{shocks} comparten ventana con otro \textit{shock} del mismo
partido, lo que motiva la decisión de no excluir las ventanas
solapadas y de absorber el cruce entre ventanas dentro de los efectos
fijos
del estimador. Una quinta parte de los \textit{shocks} tiene la
ventana pre o post recortada por el inicio o el fin del partido o por
una transición de período, condición que el estimador DiD descarta
sólo cuando afecta a la ventana post.

\begin{table}[h]
  \centering
  \small
  \caption{Composición del corpus de \textit{shocks} del torneo.}
  \label{tab:corpus-shocks}
  \begin{tabular}{llrr}
    \toprule
    Dimensión & Categoría & $n$ & \% \\
    \midrule
    Fase & Grupos        & 120 & 69{,}8 \\
         & Eliminatoria  &  52 & 30{,}2 \\
    \midrule
    Período & P1 (regulación, 1\textsuperscript{a} parte)  &  67 & 39{,}0 \\
            & P2 (regulación, 2\textsuperscript{a} parte)  & 101 & 58{,}7 \\
            & P3+P4 (prórroga)                              &   4 &  2{,}3 \\
    \midrule
    Bucket minuto & 1--15  & 16 &  9{,}3 \\
                  & 16--30 & 17 &  9{,}9 \\
                  & 31--45 & 33 & 19{,}2 \\
                  & 46--60 & 31 & 18{,}0 \\
                  & 61--75 & 31 & 18{,}0 \\
                  & 76--90 & 27 & 15{,}7 \\
                  & 91--120 & 17 &  9{,}9 \\
    \midrule
    \textit{Flags} de ventana  & Solapamiento con otro \textit{shock} & 86 & 50{,}0 \\
                               & Ventana pre truncada                  & 34 & 19{,}8 \\
                               & Sustitución en ventana                & 32 & 18{,}6 \\
                               & Ventana post truncada                 & 30 & 17{,}4 \\
                               & Tiempo extra (\textit{flag})          &  4 &  2{,}3 \\
                               & Gol en propia puerta                  &  3 &  1{,}7 \\
    \bottomrule
  \end{tabular}
\end{table}

El ATE poblacional del
\textit{shock} sobre el delta relativo al bloque del propio equipo,
estimado por DiD \textit{within-player} con errores estándar
\textit{cluster-robust} a nivel jugador, se reporta en la
tabla~\ref{tab:ate-poblacional}.

\begin{table}[h]
  \centering
  \small
  \caption{ATE del \textit{shock} sobre el delta relativo al bloque,
    por canal y dirección.}
  \label{tab:ate-poblacional}
  \begin{tabular}{llrrrrr}
    \toprule
    Canal & Contexto & ATE & SE & IC$_{95\%}$ inf. & IC$_{95\%}$ sup. & $n$ shocks \\
    \midrule
    Ofensivo   & Tras encajar &  $+0{,}0000$ & $0{,}0003$ & $-0{,}0005$ & $+0{,}0005$ & 50 \\
    Ofensivo   & Tras marcar  &  $-0{,}0001$ & $0{,}0004$ & $-0{,}0008$ & $+0{,}0007$ & 50 \\
    Defensivo  & Tras encajar &  $+0{,}0013$ & $0{,}0063$ & $-0{,}0110$ & $+0{,}0137$ & 50 \\
    Defensivo  & Tras marcar  &  $+0{,}0004$ & $0{,}0070$ & $-0{,}0133$ & $+0{,}0142$ & 50 \\
    Off-ball   & Tras encajar &  $+0{,}0000$ & $0{,}0001$ & $-0{,}0002$ & $+0{,}0002$ & 50 \\
    Off-ball   & Tras marcar  &  $-0{,}0000$ & $0{,}0001$ & $-0{,}0002$ & $+0{,}0002$ & 50 \\
    Físico     & Tras encajar &  $-0{,}0094$ & $0{,}0125$ & $-0{,}0339$ & $+0{,}0151$ & 50 \\
    Físico     & Tras marcar  &  $-0{,}0007$ & $0{,}0117$ & $-0{,}0237$ & $+0{,}0223$ & 50 \\
    \bottomrule
  \end{tabular}
\end{table}

Los ocho coeficientes son indistinguibles de cero al nivel del cinco
por ciento, con intervalos de confianza al 95\% que contienen el
origen en todas las celdas. La corrección de Benjamini-Hochberg
sobre los doce contrastes principales no produce ningún rechazo. El
\textit{shock} emocional del gol no mueve, en promedio, el rendimiento
del jugador relativo al patrón colectivo de su bloque. El análisis de potencia ex-ante sobre el
panel jugador-\textit{shock}-minuto indica que el \textit{Minimum
Detectable Effect} (MDE) efectivo al ochenta por ciento de potencia
oscila entre $0{,}0003$ desviaciones estándar para el canal
\textit{off-ball} y $0{,}036$ para el canal físico, lo que sitúa el
listón mínimo de detección de los coeficientes nulos en magnitudes
muy pequeñas. Con $N=172$ \textit{shocks} efectivos en el torneo, el
promedio poblacional no es el canal por el que se espera que aparezca
la señal \textit{clutch}.


\section{Heterogeneidad individual: el CATE jerárquico}
\label{sec:res-cate}

El ATE poblacional indistinguible de cero no implica ausencia de
efecto. La estimación del efecto heterogéneo por jugador mediante el
CATE jerárquico bayesiano \textit{multivariate} sobre los 511
jugadores con minutos en campo reporta un panorama muy distinto.
La figura \ref{fig:cate-heterogeneity} muestra la distribución del
CATE individual por canal y dirección del \textit{shock}.

\begin{figure}[h]
  \centering
  \includegraphics[width=\textwidth]{fig_cap6_cate_heterogeneity}
  \caption{Distribución del CATE individual por canal y contexto.}
  \label{fig:cate-heterogeneity}
\end{figure}

El canal ofensivo presenta la mayor varianza individual entre las
celdas: tras marcar, la desviación estándar de los CATE individuales
es $\sigma = 0{,}086$ desviaciones estándar del canal, con efectos
individuales que se extienden desde $-0{,}49$ hasta $+0{,}37$. El
canal físico también presenta varianza relevante, con
$\sigma = 0{,}032$ tras encajar. Los canales defensivo y
\textit{off-ball} presentan dispersión menor pero también
sustantivamente distinta de cero. El contraste con el ATE
poblacional reportado en la sección anterior es directo: el efecto
promedio es cero porque los efectos individuales se cancelan al
agregarlos, pero los efectos individuales en sí no son cero. La
varianza individual es la señal que el promedio esconde, y es la
cantidad sobre la que opera la lectura \textit{scout-facing} del
trabajo.


\section{Jugadores con respuesta significativa por contexto}
\label{sec:res-scout-facing}

Antes de entrar en los nombres concretos, la figura~\ref{fig:scatter-remontador-cerrojo}
ofrece el panorama del torneo sobre los dos índices agregados que el
\CCV reporta: el Remontador (eje horizontal, respuesta tras encajar)
y el Cerrojo (eje vertical, respuesta tras marcar). Cada punto es un
jugador del Mundial Qatar 2022 con muestra suficiente, coloreado por
su \textit{bucket} posicional. El gráfico permite leer perfiles
extremos en cada cuadrante ---el doble-clutch arriba a la derecha, el
doble-choke abajo a la izquierda--- y comprobar visualmente que la
mayoría de jugadores se concentra alrededor del origen, en la zona
donde el efecto individual está dentro del ruido del estimador.

\begin{figure}[h]
  \centering
  \includegraphics[width=\textwidth]{scatter_remontador_cerrojo}
  \caption{Panorama del torneo: Remontador frente a Cerrojo.}
  \label{fig:scatter-remontador-cerrojo}
\end{figure}

La tabla~\ref{tab:sig-por-celda} resume el conteo de jugadores
significativos en cada una de las ocho celdas canal-por-contexto y en
los tres umbrales del índice \textit{pressure-clutch}. La señal
individual se concentra en dos celdas: el canal ofensivo tras marcar
(22 jugadores) y la respuesta a la proximidad de eliminación (entre
4 y 10 según el umbral). Las seis celdas restantes son
\textit{inconclusive} para todos los jugadores del corpus, lo que
encaja con la heterogeneidad reportada en la sección anterior: las
celdas con $\sigma$ del CATE individual más pequeño concentran poca
señal por jugador y no producen significatividad posterior individual
con $N=172$ \textit{shocks}. La
discusión operativa de esta asimetría se retoma en el
capítulo~\ref{ch:discusion}.

\begin{table}[h]
  \centering
  \small
  \caption{Jugadores significativos por celda y por umbral.}
  \label{tab:sig-por-celda}
  \begin{tabular}{llrrr}
    \toprule
    Canal & Contexto & PRO & ANTI & Total \\
    \midrule
    Ofensivo   & Tras encajar &  0 &  0 &  0 \\
    Ofensivo   & Tras marcar  &  5 & 17 & 22 \\
    Defensivo  & Tras encajar &  0 &  0 &  0 \\
    Defensivo  & Tras marcar  &  0 &  0 &  0 \\
    Off-ball   & Tras encajar &  0 &  0 &  0 \\
    Off-ball   & Tras marcar  &  0 &  0 &  0 \\
    Físico     & Tras encajar &  0 &  0 &  0 \\
    Físico     & Tras marcar  &  0 &  0 &  0 \\
    \midrule
    \multicolumn{2}{l}{\textit{Pressure-clutch} (umbral $0{,}80$)}      &  6 &  4 & 10 \\
    \multicolumn{2}{l}{\textit{Pressure-clutch} (umbral $0{,}85$)}      &  3 &  1 &  4 \\
    \multicolumn{2}{l}{\textit{Pressure-clutch} (umbral $0{,}95$)}      &  1 &  0 &  1 \\
    \bottomrule
  \end{tabular}
\end{table}

Los jugadores con efecto significativo en el contexto de proximidad
de eliminación alta ---el indicador \textit{pressure-clutch} del
\CCV--- se reportan en la tabla~\ref{tab:pressure-clutch}. La
significatividad se define sobre la probabilidad posterior
$P(\beta > 0 \mid \text{datos}) \geq 0{,}85$ para el efecto positivo
o $P(\beta > 0 \mid \text{datos}) \leq 0{,}15$ para el efecto
negativo.

\begin{table}[h]
  \centering
  \small
  \caption{Jugadores con efecto \textit{pressure-clutch} significativo.}
  \label{tab:pressure-clutch}
  \begin{tabular}{llrrrl}
    \toprule
    Jugador & Selección & Media post. & IC$_{80\%}$ inf. & IC$_{80\%}$ sup. & $P(\beta>0)$ \\
    \midrule
    Kylian Mbappé   & Francia    & $+0{,}110$ & $+0{,}029$ & $+0{,}193$ & $0{,}97$ \\
    Marcus Rashford & Inglaterra & $+0{,}050$ & $-0{,}003$ & $+0{,}110$ & $0{,}89$ \\
    Bukayo Saka     & Inglaterra & $+0{,}051$ & $-0{,}009$ & $+0{,}116$ & $0{,}86$ \\
    \midrule
    Mohammed Muntari & Qatar     & $-0{,}055$ & $-0{,}120$ & $+0{,}004$ & $0{,}12$ \\
    \bottomrule
  \end{tabular}
\end{table}

La figura~\ref{fig:radar-mbappe} reporta la ficha
\textit{scout-facing} de Mbappé como ilustración del formato por
jugador del \CCV: doce ejes que cruzan los cuatro canales con los
tres contextos del \textit{shock}, expresados como percentil dentro
de su posición sobre el corpus del torneo. La lectura del radar
requiere una precisión metodológica relevante: el percentil que
muestra cada eje no es una medida de significatividad del efecto
causal sobre ese jugador, sino su \textit{ranking} relativo dentro
de su \textit{bucket} posicional. Un percentil alto indica que el
jugador está entre los más altos de su posición en esa celda, pero
la magnitud absoluta del efecto puede seguir siendo indistinguible
de cero si la probabilidad posterior por jugador no cruza el umbral
de significatividad. Las dos lecturas ---percentil posicional y
probabilidad posterior por jugador--- son complementarias y se
reportan conjuntamente en la tabla canónica del \CCV; el anexo
\ref{anexo:radares} reproduce las fichas equivalentes para Messi,
Hakimi y Brozović.

\begin{figure}[h]
  \centering
  \includegraphics[width=0.95\textwidth]{radar_report_3870}
  \caption{Ficha \textit{scout-facing} de Kylian Mbappé.}
  \label{fig:radar-mbappe}
\end{figure}

Los tres jugadores con efecto positivo significativo están liderados
por Kylian Mbappé, con un coeficiente de $+0{,}110$ desviaciones
estándar y probabilidad posterior $0{,}97$. La identificación es
consistente con la observación cualitativa de su irrupción en el
cierre de la final Argentina-Francia, donde concentró cuatro de los
seis goles del partido en una ventana corta de máxima presión por
eliminación. Marcus Rashford y Bukayo Saka, ambos por Inglaterra,
completan el grupo positivo con magnitudes equivalentes en torno a
$+0{,}050$ desviaciones estándar y probabilidades posteriores
respectivamente $0{,}89$ y $0{,}86$. El caso negativo significativo
corresponde a Mohammed Muntari, cuyo coeficiente cae por debajo de
cero con probabilidad posterior $P(\beta < 0) = 0{,}88$ ---es decir,
$P(\beta > 0) = 0{,}12$ en la tabla~\ref{tab:pressure-clutch}.

A nivel canal-específico, la celda ofensivo-tras-marcar concentra la
mayor cantidad de efectos individuales significativos: veintidós
jugadores presentan respuesta significativa al \textit{shock} de
marcar sobre su producción ofensiva inmediata, cinco con coeficiente
positivo y diecisiete con coeficiente negativo. La
figura~\ref{fig:scatter-ataque-marcar-encajar} reporta el panorama
del torneo sobre las dos celdas ofensivas: ataque tras marcar (eje
horizontal) y ataque tras encajar (eje vertical), con un punto por
jugador.

\begin{figure}[h]
  \centering
  \includegraphics[width=\textwidth]{scatter_ataque_marcar_encajar}
  \caption{Panorama del torneo: ataque tras marcar frente a ataque
    tras encajar.}
  \label{fig:scatter-ataque-marcar-encajar}
\end{figure}

Los cinco con efecto positivo son Eric Maxim Choupo-Moting, Dušan
Vlahović, Raheem Sterling, Harry Kane y Cristiano Ronaldo, en orden
decreciente de probabilidad posterior. La narrativa de cada caso ---celebración
seguida de involucración mayor en jugadas posteriores, búsqueda
activa de un segundo gol--- es consistente con los grades humanos
PFF para esos mismos jugadores en los partidos del torneo. Los
diecisiete con efecto negativo incluyen al propio Mbappé, cuyo CATE
ofensivo-tras-marcar es significativamente negativo
($P(\beta > 0) = 0{,}004$): tras anotar, su producción ofensiva
inmediata cae respecto al patrón colectivo de Francia. La
combinación es coherente: Mbappé responde fuertemente bajo presión
de eliminación pero relaja su intensidad ofensiva en la ventana
inmediata después de marcar, dejando momentáneamente que el bloque
recupere el control del partido.

En los agregados de la tabla \textit{scout-facing} ---el Remontador
post-encajar y el Cerrojo post-marcar--- no se identifican jugadores
significativos a nivel poblacional. El motivo es matemático: cada
agregado es la suma de cuatro celdas canal-por-contexto cuyos
efectos individuales tienden a no estar correlacionados en signo,
lo que reduce la señal individual al promediarlos. La respuesta
condicional al contexto de eliminación inminente, en cambio, es
identificable porque opera sobre un eje único de modulación.


\section{Validación cualitativa contra los \textit{grades} humanos PFF}
\label{sec:res-validacion-pff}

Como verificación externa cualitativa, se calcula la correlación de
Spearman entre las puntuaciones del \CCV por canal y por índice
agregado y el \textit{grade} humano PFF agregado por jugador, sobre
los 166 jugadores con al menos 270\,minutos jugados y muestra fiable.
La cifra es modesta en términos absolutos ---un tercio del corpus
total del torneo--- y conviene tenerlo presente al interpretar los
coeficientes: la validación direccional opera sobre una muestra
acotada, no sobre una población amplia de \textit{labels} humanos. La
tabla~\ref{tab:correlacion-pff} reporta los coeficientes.

\begin{table}[h]
  \centering
  \small
  \caption{Correlación de Spearman entre \CCV y \textit{grade}
    humano PFF.}
  \label{tab:correlacion-pff}
  \begin{tabular}{lrrl}
    \toprule
    Métrica \CCV & $\rho$ Spearman & $p$ & Significativa \\
    \midrule
    Canal ofensivo (\texttt{ccv\_atk})         & $+0{,}017$ & $0{,}83$  & --- \\
    Canal defensivo (\texttt{ccv\_def})        & $-0{,}071$ & $0{,}37$  & --- \\
    Canal \textit{off-ball} (\texttt{ccv\_off}) & $+0{,}211$ & $0{,}006$ & 1\% \\
    Canal físico (\texttt{ccv\_phys})          & $+0{,}090$ & $0{,}25$  & --- \\
    \midrule
    Índice Remontador (chasing)         & $+0{,}177$ & $0{,}022$ & 5\% \\
    Índice Cerrojo (protecting)         & $+0{,}005$ & $0{,}94$  & --- \\
    Índice \textit{Pressure Response}   & $-0{,}122$ & $0{,}12$  & --- \\
    \bottomrule
  \end{tabular}
\end{table}

El patrón es modesto y heterogéneo. El canal \textit{off-ball} es el
único que correlaciona significativamente al uno por ciento con el
\textit{grade} humano agregado ($\rho = +0{,}211$); el índice
Remontador lo hace al cinco por ciento ($\rho = +0{,}177$). Los demás
coeficientes no se distinguen de cero con la muestra disponible.
El \textit{grade} humano PFF disponible es un valor único agregado
por jugador y partido que sintetiza el rendimiento global, no una
descomposición por canal o por contexto de \textit{shock}. Los
canales del \CCV, en cambio, son métricas \textit{event-level}
agregadas por minuto y residualizadas contra el bloque colectivo del
propio equipo en cada \textit{shock}. Las dos medidas cubren
propiedades parcialmente solapadas pero no idénticas, y eso explica
que la correlación varíe tanto entre canales. La validación externa
no es contundente: valida direccionalmente que el canal
\textit{off-ball} y el agregado Remontador recogen propiedades
reconocibles por evaluadores humanos expertos, sin pretender que los
CATE y los \textit{grades} midan lo mismo.
